\documentclass[12pt, a4paper]{article}

% ---------------------------------------------------------------
% Packages
% ---------------------------------------------------------------
\usepackage[T1]{fontenc}
\usepackage[utf8]{inputenc}
\usepackage{amsmath, amssymb, amsthm}
\usepackage{mathtools}
\usepackage{booktabs}
\usepackage{array}
\usepackage{xcolor}
\usepackage{hyperref}
\usepackage{geometry}
\usepackage{enumitem}
\usepackage{tcolorbox}
\usepackage{mdframed}

\geometry{margin=2.5cm}

\hypersetup{
  colorlinks = true,
  linkcolor  = blue!60!black,
  citecolor  = blue!60!black,
}

% ---------------------------------------------------------------
% Custom environments
% ---------------------------------------------------------------
\newtheoremstyle{condstyle}{}{}{\itshape}{}{\bfseries}{.}{.5em}{}
\theoremstyle{condstyle}
\newtheorem{condition}{Condition}[section]

\theoremstyle{plain}
\newtheorem{proposition}{Proposition}[section]
\newtheorem{corollary}{Corollary}[section]
\newtheorem{remark}{Remark}[section]

% Shaded box for key results
\tcbuselibrary{skins}
\newtcolorbox{keyresult}[1][]{
  enhanced,
  colback=blue!4!white,
  colframe=blue!40!black,
  fonttitle=\bfseries,
  title=#1,
  left=6pt, right=6pt, top=4pt, bottom=4pt,
  arc=3pt,
}

\newtcolorbox{notebox}[1][]{
  enhanced,
  colback=orange!5!white,
  colframe=orange!60!black,
  fonttitle=\bfseries,
  title=#1,
  left=6pt, right=6pt, top=4pt, bottom=4pt,
  arc=3pt,
}

% ---------------------------------------------------------------
% Notation macros
% ---------------------------------------------------------------
\newcommand{\EN}{\mathbb{E}_N}
\newcommand{\bD}{\mathbf{D}}
\newcommand{\bY}{\mathbf{Y}}
\newcommand{\bZ}{\mathbf{Z}}
\newcommand{\bone}{\mathbf{1}_N}
\newcommand{\bzero}{\mathbf{0}_N}
\newcommand{\IN}{\mathbf{I}_N}
\newcommand{\tY}{\tilde{Y}}
\newcommand{\tD}{\tilde{D}}
\newcommand{\tZ}{\tilde{Z}}
\newcommand{\that}[1]{\hat{\tau}_{#1}}
\newcommand{\phat}{\hat{p}}
\newcommand{\kone}{\kappa_1}
\newcommand{\kzero}{\kappa_0}

% ---------------------------------------------------------------
% Document
% ---------------------------------------------------------------
\begin{document}

\begin{titlepage}
\centering
\vspace*{2cm}
{\LARGE\bfseries Outcome Weights for Kappa-Based LATE Estimators}\\[0.5em]
{\large A Step-by-Step Derivation in the PIVE Framework of Knaus (2024)}\\[2cm]

{\large Working Notes}\\[0.5em]
{\normalsize Based on Knaus (2024) and Słoczy\'nski, Uysal \& Wooldridge (2024)}\\[3cm]

\begin{abstract}
\noindent
These notes derive the outcome weights $\omega_i$ such that $\hat\tau = \sum_i \omega_i Y_i$
for three kappa-based estimators of the Local Average Treatment Effect (LATE):
the fully-normalized estimator $\that{a,10}$, the unnormalized Abadie estimator $\that{a}$,
and the Uysal normalized estimator $\that{u}$.
The derivation follows the Pseudo-IV Estimator (PIVE) framework of Knaus (2024).
For each estimator we state which conditions from Knaus (2024) are needed,
explain what those conditions mean in the kappa context,
and show step by step whether the outcome weights are
fully-normalized ($\sum_i \omega_i D_i = 1$, $\sum_i \omega_i (1-D_i) = -1$),
scale-normalized, or unnormalized.
\end{abstract}
\vfill
\end{titlepage}

\tableofcontents
\newpage

% ===============================================================
\section{Setting and Notation}
\label{sec:setup}
% ===============================================================

We work in a binary instrumental variables (IV) setting.
The observed data for unit $i = 1, \ldots, N$ are
\[
  O_i = (D_i,\; X_i,\; Y_i,\; Z_i)',
\]
where $D_i \in \{0,1\}$ is a binary treatment, $Z_i \in \{0,1\}$ is a binary instrument,
$Y_i$ is the outcome of interest, and $X_i$ is a vector of pre-treatment covariates.

\medskip
\noindent\textbf{Target parameter.}
The target is the Local Average Treatment Effect (LATE),
\[
  \tau_{\mathrm{LATE}} = \mathbb{E}[Y_i(1) - Y_i(0) \mid \text{Complier}_i],
\]
where a \emph{complier} is a unit whose treatment status is shifted by the instrument.

\medskip
\noindent\textbf{Instrument propensity score.}
Let $p(X_i) = \mathbb{E}[Z_i \mid X_i]$ denote the true instrument propensity score,
and $\phat_i = \hat{p}(X_i)$ its estimator (typically logistic regression or CBPS).

\medskip
\noindent\textbf{Matrix notation.}
Bold letters denote $N$-vectors or $N\times N$ matrices.
$\bone$ and $\bzero$ are vectors of ones and zeros of length $N$.
$\mathbf{D} = (D_1,\ldots,D_N)'$ and $\mathbf{Y} = (Y_1,\ldots,Y_N)'$.

% ===============================================================
\section{The Three Kappa Weights}
\label{sec:kappa}
% ===============================================================

The entire family of estimators studied in S\l oczy\'nski et al.\ (2024) is built from three scalar weights
per observation, first introduced by Abadie (2003):

\begin{align}
  \kappa_i   &= 1 - D_i\,\frac{1-Z_i}{1-\phat_i} - (1-D_i)\,\frac{Z_i}{\phat_i},
  \label{eq:kappa}\\[6pt]
  \kappa_{i1} &= D_i\,\frac{Z_i - \phat_i}{\phat_i(1-\phat_i)},
  \label{eq:kappa1}\\[6pt]
  \kappa_{i0} &= (1-D_i)\,\frac{(1-Z_i)-(1-\phat_i)}{\phat_i(1-\phat_i)}
               = -(1-D_i)\,\frac{Z_i - \phat_i}{\phat_i(1-\phat_i)}.
  \label{eq:kappa0}
\end{align}

\begin{notebox}[Intuition]
$\kappa_{i1}$ is nonzero \emph{only} for treated units ($D_i=1$).
It is positive when $Z_i=1$ and negative when $Z_i=0$,
isolating the complier contribution from the treated group.
$\kappa_{i0}$ works symmetrically for untreated units ($D_i=0$).
Together, $\kappa_1$ and $\kappa_0$ identify the complier subpopulation from both directions.
Note also that $\kappa_{i1} \cdot (1-D_i) = 0$ and $\kappa_{i0} \cdot D_i = 0$ for all $i$
--- a fact we will use extensively in the normalization proofs.
\end{notebox}

In vector notation:
\[
  \boldsymbol{\kappa} = (\kappa_1,\ldots,\kappa_N)',\quad
  \boldsymbol{\kappa}_1 = (\kappa_{11},\ldots,\kappa_{N1})',\quad
  \boldsymbol{\kappa}_0 = (\kappa_{10},\ldots,\kappa_{N0})'.
\]

% ===============================================================
\section{The PIVE Framework}
\label{sec:pive}
% ===============================================================

Knaus (2024) defines a general class of \emph{Pseudo-IV Estimators} (PIVEs)
as estimators solving an empirical moment condition of the form
\begin{equation}
  \EN\!\left[(\tY_i - \hat\tau\,\tD_i)\,\tZ_i\right] = 0,
  \label{eq:moment}
\end{equation}
with scalar pseudo-outcome $\tY_i = f_Y(O_i;\hat\eta^Y_i)$,
pseudo-treatment $\tD_i = f_D(O_i;\hat\eta^D_i)$,
and pseudo-instrument $\tZ_i = f_Z(O_i;\hat\eta^Z_i)$.
Solving \eqref{eq:moment} yields
\begin{equation}
  \hat\tau = \frac{\EN[\tZ_i\tY_i]}{\EN[\tZ_i\tD_i]}
            = (\tilde{\mathbf{Z}}'\tilde{\mathbf{D}})^{-1}\tilde{\mathbf{Z}}'\tilde{\mathbf{Y}}.
  \label{eq:pive_sol}
\end{equation}

\begin{proposition}[Outcome weights of PIVEs, Knaus 2024, Prop.~1]
\label{prop:weights}
If $\tilde{\mathbf{Z}}'\tilde{\mathbf{D}} \neq 0$ and there exists a unique $N\times N$
transformation matrix $\mathbf{T}$ such that $\mathbf{T}\mathbf{Y} = \tilde{\mathbf{Y}}$,
then the outcome weights of a PIVE take the closed form
\begin{equation}
  \boldsymbol{\omega}' = (\tilde{\mathbf{Z}}'\tilde{\mathbf{D}})^{-1}\,\tilde{\mathbf{Z}}'\,\mathbf{T},
  \label{eq:weights_general}
\end{equation}
so that $\hat\tau = \boldsymbol{\omega}'\mathbf{Y} = \sum_{i=1}^N \omega_i Y_i$.
\end{proposition}

\medskip
\noindent\textbf{Shortcuts for normalization.}
To classify the resulting outcome weights,
Knaus (2024) provides three matrix shortcuts (Section~4.1):

\begin{enumerate}[label=\textbf{SC\arabic*.}]
  \item $\mathbf{T}\bone = \bzero$ implies $\sum_i \omega_i = 0$
        \hfill(necessary for normalization)
  \item $\mathbf{T}\mathbf{D} = \tilde{\mathbf{D}}$ implies $\sum_i \omega_i D_i = 1$
        \hfill(treated weights sum to one)
  \item $\mathbf{T}(\bone - \mathbf{D}) = -\tilde{\mathbf{D}}$ implies $\sum_i \omega_i(1-D_i) = -1$
        \hfill(untreated weights sum to minus one)
\end{enumerate}

\noindent Full normalization ($\sum_i\omega_i=0$, $\sum_i\omega_iD_i=1$, $\sum_i\omega_i(1-D_i)=-1$)
follows from SC1 + SC2 (since SC3 is then implied by subtraction).

\medskip
\noindent\textbf{Important caveat for the kappa estimators.}
Knaus (2024) notes that for $\that{a,10}$,
the shortcuts SC1--SC3 via $\mathbf{T}$ alone are \emph{not sufficient}
because the normalizing scalars $({\bone}'\boldsymbol{\kappa}_1)^{-1}$ and
$({\bone}'\boldsymbol{\kappa}_0)^{-1}$ are absorbed into the pseudo-instrument
rather than $\mathbf{T}$.
We must therefore check the \emph{full weight expression} directly
(see Appendix~A.4 of Knaus 2024 and Section~\ref{sec:a10} below).

% ===============================================================
\section{Which Conditions Does Knaus (2024) Introduce?}
\label{sec:conditions}
% ===============================================================

Knaus (2024) introduces six conditions in Section~4.2.1 that govern
when outcome weights exist and what normalization properties they possess.
Below we state all six, explain their general content, and note
which ones are relevant for the kappa estimators in our context.

\begin{condition}[Smoother matrix]\label{cond:smoother}
  A unique smoother matrix exists that produces the outcome nuisance predictions
  by left-multiplying the outcome vector:
  \begin{align*}
    \text{(C1a)} &\quad \mathbf{S}\mathbf{Y} = \hat{\mathbf{Y}},\\
    \text{(C1b)} &\quad \mathbf{S}^d_0\mathbf{Y} = \hat{\mathbf{Y}}^{d_0},\quad
                         \mathbf{S}^d_1\mathbf{Y} = \hat{\mathbf{Y}}^{d_1},\\
    \text{(C1c)} &\quad \mathbf{S}^z_0\mathbf{Y} = \hat{\mathbf{Y}}^{z_0},\quad
                         \mathbf{S}^z_1\mathbf{Y} = \hat{\mathbf{Y}}^{z_1}.
  \end{align*}
  \textit{Purpose:} This is the existence condition for outcome weights.
  Without a smoother matrix, the pseudo-outcome $\tY$ cannot be written as $\mathbf{T}\mathbf{Y}$,
  so Proposition~\ref{prop:weights} does not apply.\\
  \textit{In our context:} The kappa estimators $\that{a}$, $\that{a,10}$, and $\that{u}$
  involve \emph{no outcome regression whatsoever} --- only $Z_i$, $D_i$, and $\phat_i$.
  Condition~\ref{cond:smoother} is therefore \textbf{automatically satisfied} (trivially,
  with the identity matrix as smoother), and outcome weights exist unconditionally.
\end{condition}

\begin{condition}[Covariate matrix with constant]\label{cond:constant}
  The covariate matrix $\mathbf{X}$ contains a column of ones as its first column,
  so that $\mathbf{X}(1, \mathbf{0}_p')' = \bone$.\\
  \textit{Purpose:} Ensures OLS and TSLS are normalized.\\
  \textit{In our context:} Not relevant for kappa estimators.
\end{condition}

\begin{condition}[Affine smoother matrix]\label{cond:affine}
  All rows of each smoother matrix sum to one:
  $\mathbf{S}\bone = \mathbf{S}^d_0\bone = \mathbf{S}^d_1\bone =
   \mathbf{S}^z_0\bone = \mathbf{S}^z_1\bone = \bone$.\\
  \textit{Purpose:} Ensures normalized weights for AIPW, PLR, and causal/instrumental forests.
  Most regression smoothers (random forests, kernel regression, ridge) satisfy this;
  boosted trees may not.\\
  \textit{In our context:} Not directly required for kappa estimators,
  but the analogous role is played by the explicit normalization built into
  $\that{a,10}$ and $\that{u}$.
\end{condition}

\begin{condition}[No smoothing between treatment groups]\label{cond:nosmoothbetween}
  Predictions $\hat{Y}^{d_1}$ are formed using treated units only,
  and $\hat{Y}^{d_0}$ using untreated units only:
  $\mathbf{S}^d_1(\bone - \mathbf{D}) = \bzero$ and $\mathbf{S}^d_0\mathbf{D} = \bzero$.\\
  \textit{Purpose:} Ensures AIPW is fully normalized.\\
  \textit{In our context:} Not relevant for kappa estimators.
\end{condition}

\begin{condition}[Outcome smoother applied to treatment]\label{cond:same_smoother}
  \begin{align*}
    \text{(C5a)} &\quad \mathbf{S}\mathbf{D} = \hat{\mathbf{D}},\\
    \text{(C5b)} &\quad \mathbf{S}^z_1\mathbf{D} = \hat{\mathbf{D}}^{z_1},\quad
                          \mathbf{S}^z_0\mathbf{D} = \hat{\mathbf{D}}^{z_0}.
  \end{align*}
  \textit{Purpose:} Ensures full normalization for PLR, causal/instrumental forests (C5a),
  and Wald-AIPW (C5b).
  This requires using the \emph{same} model to predict both outcome and treatment,
  which conflicts with standard flexible implementations.\\
  \textit{In our context:} Not relevant for kappa estimators.
\end{condition}

\begin{condition}[Normalized IPW weights]\label{cond:ipwnorm}
  \begin{align*}
    \text{(C6a)} &\quad \lambda^{\mathrm{norm}}_{1,i} = \lambda^{\mathrm{ipw}}_{1,i}\big/\EN[\lambda^{\mathrm{ipw}}_1],\quad
                         \bone'\boldsymbol{\lambda}^{\mathrm{norm}}_1 = \bone'\boldsymbol{\lambda}^{\mathrm{norm}}_0 = N,\\
    \text{(C6b)} &\quad \text{analogous normalization for instrument-group IPW weights.}
  \end{align*}
  \textit{Purpose:} Ensures full normalization for standard IPW (C6a) and Wald-IPW (C6b).
  This is the Hájek (1971) normalization.\\
  \textit{In our context:} $\that{u}$ is the direct analogue of Wald-IPW with C6b applied.
  The separate normalization in $\that{a,10}$ is a different mechanism that achieves
  the same effect without explicitly invoking C6b.
\end{condition}

\begin{notebox}[Summary for kappa estimators]
Among the six conditions, only the trivial case of Condition~\ref{cond:smoother} matters
for the kappa estimators. No smoother for outcome regressions is needed,
so outcome weights always exist. Normalization properties then follow purely from
the algebraic structure of the kappa weights themselves.
Conditions \ref{cond:constant}--\ref{cond:same_smoother} are irrelevant.
Condition~\ref{cond:ipwnorm} is implicitly used by $\that{u}$.
\end{notebox}

% ===============================================================
\section{The Five Estimators}
\label{sec:estimators}
% ===============================================================

S\l oczy\'nski et al.\ (2024) study five estimators. In vector form, using the kappa definitions
from Section~\ref{sec:kappa}:

\begin{align}
  \that{a}    &= (\bone'\boldsymbol{\kappa})^{-1}\,\bone'\,\mathrm{diag}(\boldsymbol{\kappa}_1 - \boldsymbol{\kappa}_0)\,\mathbf{Y},
  \label{eq:tau_a}\\[4pt]
  \that{a,1}  &= (\bone'\boldsymbol{\kappa}_1)^{-1}\,\bone'\,\mathrm{diag}(\boldsymbol{\kappa}_1 - \boldsymbol{\kappa}_0)\,\mathbf{Y},
  \label{eq:tau_a1}\\[4pt]
  \that{a,0}  &= (\bone'\boldsymbol{\kappa}_0)^{-1}\,\bone'\,\mathrm{diag}(\boldsymbol{\kappa}_1 - \boldsymbol{\kappa}_0)\,\mathbf{Y},
  \label{eq:tau_a0}\\[4pt]
  \that{a,10} &= N^{-1}\,\bone'\,\mathrm{diag}\!\left[\boldsymbol{\kappa}_1(\bone'\boldsymbol{\kappa}_1)^{-1}N
                  - \boldsymbol{\kappa}_0(\bone'\boldsymbol{\kappa}_0)^{-1}N\right]\mathbf{Y},
  \label{eq:tau_a10}\\[4pt]
  \that{u}    &= \frac{%
    \displaystyle\left(\sum_i \frac{Z_i}{\phat_i}\right)^{-1}\!\sum_i \frac{Z_i Y_i}{\phat_i}
    - \left(\sum_i \frac{1-Z_i}{1-\phat_i}\right)^{-1}\!\sum_i \frac{(1-Z_i)Y_i}{1-\phat_i}
  }{%
    \displaystyle\left(\sum_i \frac{Z_i}{\phat_i}\right)^{-1}\!\sum_i \frac{Z_i D_i}{\phat_i}
    - \left(\sum_i \frac{1-Z_i}{1-\phat_i}\right)^{-1}\!\sum_i \frac{(1-Z_i)D_i}{1-\phat_i}
  }.
  \label{eq:tau_u}
\end{align}

\noindent
Estimators \eqref{eq:tau_a}--\eqref{eq:tau_a0} share the same numerator
$\bone'\,\mathrm{diag}(\boldsymbol{\kappa}_1 - \boldsymbol{\kappa}_0)\,\mathbf{Y}$
but use different denominators. In scalar form:
\[
  \that{a}   = \frac{\bar\kappa_1 - \bar\kappa_0}{\bar\kappa},\quad
  \that{a,1} = \frac{\bar\kappa_1 - \bar\kappa_0}{\bar\kappa_1},\quad
  \that{a,0} = \frac{\bar\kappa_1 - \bar\kappa_0}{\bar\kappa_0},
\]
where $\bar\kappa = N^{-1}\sum_i \kappa_i Y_i$ and so on.

% ===============================================================
\section{Derivation: \texorpdfstring{$\hat\tau_{a,10}$}{tau\_a,10}}
\label{sec:a10}
% ===============================================================

\subsection{PIVE representation}

Write \eqref{eq:tau_a10} compactly as
\begin{equation}
  \that{a,10}
  = N^{-1}\,\bone'\underbrace{%
    \mathrm{diag}\!\left[\boldsymbol{\kappa}_1(\bone'\boldsymbol{\kappa}_1)^{-1}N
                        - \boldsymbol{\kappa}_0(\bone'\boldsymbol{\kappa}_0)^{-1}N\right]
    }_{=:\mathbf{T}^{a,10}}\mathbf{Y}.
  \label{eq:ta10_compact}
\end{equation}

Matching to the PIVE structure \eqref{eq:pive_sol}:
\begin{itemize}
  \item $\tilde{\mathbf{Z}}' = N^{-1}\bone'$ \quad (pseudo-instrument is constant $1/N$ per observation),
  \item $\tilde{\mathbf{D}} = \bone$ \quad (pseudo-treatment is identically one),
  \item $\tilde{\mathbf{Z}}'\tilde{\mathbf{D}} = N^{-1}\bone'\bone = 1$ \quad (normalizing scalar = 1),
  \item $\mathbf{T}^{a,10} = \mathrm{diag}\!\left[\boldsymbol{\kappa}_1(\bone'\boldsymbol{\kappa}_1)^{-1}N
          - \boldsymbol{\kappa}_0(\bone'\boldsymbol{\kappa}_0)^{-1}N\right]$.
\end{itemize}

\noindent Since the kappa weights are purely functions of $Z_i$, $D_i$, $\phat_i$
and \emph{not} of $Y_i$, the smoother matrix is trivially available
(Condition~\ref{cond:smoother} holds without restriction).
By Proposition~\ref{prop:weights}, the outcome weights are:
\begin{equation}
  \boldsymbol{\omega}^{a,10\prime}
  = (\tilde{\mathbf{Z}}'\tilde{\mathbf{D}})^{-1}\,\tilde{\mathbf{Z}}'\,\mathbf{T}^{a,10}
  = N^{-1}\,\bone'\,\mathrm{diag}\!\left[
      \boldsymbol{\kappa}_1(\bone'\boldsymbol{\kappa}_1)^{-1}N
      - \boldsymbol{\kappa}_0(\bone'\boldsymbol{\kappa}_0)^{-1}N
    \right].
  \label{eq:w_a10}
\end{equation}

\noindent The per-observation outcome weight is therefore
\begin{equation}
  \boxed{
    \omega_i^{a,10}
    = \frac{\kappa_{i1}}{\sum_j \kappa_{j1}}
    - \frac{\kappa_{i0}}{\sum_j \kappa_{j0}}.
  }
  \label{eq:w_a10_scalar}
\end{equation}

\begin{notebox}[Interpretation]
Each observation's weight is its share of the treated-side kappa
minus its share of the untreated-side kappa.
Treated units ($D_i=1$) receive positive weight through $\kappa_{i1}$;
untreated units ($D_i=0$) receive negative weight through $-\kappa_{i0}$.
Units with $D_i=1$ and $Z_i=0$ (defiers in a no-defiers world, or always-takers)
get negative weight via $\kappa_{i1}<0$.
\end{notebox}

\subsection{Normalization: weights sum to zero}

We check normalization by computing $\sum_i \omega_i^{a,10}$ directly from \eqref{eq:w_a10}
(the transformation-matrix shortcut SC1 is not sufficient here because the
normalizing scalars sit outside $\mathbf{T}$, as noted in Section~\ref{sec:pive}).

\begin{align*}
  \sum_i \omega_i^{a,10}
  &= N^{-1}\,\bone'\,\mathbf{T}^{a,10}\,\bone\\
  &= N^{-1}\,\bone'\!\left[\boldsymbol{\kappa}_1(\bone'\boldsymbol{\kappa}_1)^{-1}N
                            - \boldsymbol{\kappa}_0(\bone'\boldsymbol{\kappa}_0)^{-1}N\right]\\
  &= N^{-1}\left[N\,(\bone'\boldsymbol{\kappa}_1)(\bone'\boldsymbol{\kappa}_1)^{-1}
               - N\,(\bone'\boldsymbol{\kappa}_0)(\bone'\boldsymbol{\kappa}_0)^{-1}\right]\\
  &= N^{-1}\left[N - N\right] = 0. \qquad \checkmark
\end{align*}

\begin{keyresult}[$\that{a,10}$ is normalized (no conditions required)]
$\displaystyle\sum_{i=1}^N \omega_i^{a,10} = 0$.
This follows unconditionally from the built-in normalization of $\boldsymbol{\kappa}_1$
and $\boldsymbol{\kappa}_0$ in \eqref{eq:tau_a10}.
No smoother condition (C1--C3) or IPW normalization condition (C6) is needed.
\end{keyresult}

\subsection{Full normalization: treated weights sum to one}

We compute $\sum_i \omega_i^{a,10} D_i$ directly:

\begin{align*}
  \sum_i \omega_i^{a,10} D_i
  &= N^{-1}\,\bone'\,\mathbf{T}^{a,10}\,\mathbf{D}\\
  &= N^{-1}\,\bone'\!\left[\boldsymbol{\kappa}_1(\bone'\boldsymbol{\kappa}_1)^{-1}N
                            - \boldsymbol{\kappa}_0(\bone'\boldsymbol{\kappa}_0)^{-1}N\right]\circ\mathbf{D}
  \intertext{where $\circ$ denotes element-wise multiplication.
  Since $\kappa_{i0}(D_i=1)=0$ for all $i$ (untreated-side kappa is zero for treated units),
  and $\kappa_{i1}(D_i=1) = \kappa_{i1}$ (treated-side kappa is zero for untreated, so multiplying by $D$ keeps it):}
  &= N^{-1}\,\left[N\,(\bone'\boldsymbol{\kappa}_1)^{-1}(\bone'\boldsymbol{\kappa}_1)
                   - 0\right]\\
  &= N^{-1} \cdot N = 1. \qquad\checkmark
\end{align*}

\begin{keyresult}[$\that{a,10}$ is fully normalized (no conditions required)]
$\displaystyle\sum_{i=1}^N \omega_i^{a,10} D_i = 1$.
Combined with $\sum_i \omega_i^{a,10} = 0$, it follows by subtraction that
$\displaystyle\sum_{i=1}^N \omega_i^{a,10}(1-D_i) = -1$.
\textbf{Full normalization holds unconditionally} --- no conditions from Knaus (2024)
need to be invoked.
This stands in contrast to PLR, causal forests, and Wald-AIPW,
which require Conditions~\ref{cond:affine}~and~\ref{cond:same_smoother} to achieve
the same property.
\end{keyresult}

% ===============================================================
\section{Derivation: \texorpdfstring{$\hat\tau_a$}{tau\_a} (Fully Unnormalized)}
\label{sec:tau_a}
% ===============================================================

\subsection{PIVE representation}

From \eqref{eq:tau_a}, the estimator in PIVE form is:
\begin{equation}
  \that{a}
  = (\bone'\boldsymbol{\kappa})^{-1}\,\bone'\,
    \underbrace{\mathrm{diag}(\boldsymbol{\kappa}_1 - \boldsymbol{\kappa}_0)}_{=:\mathbf{T}^{a}}\,\mathbf{Y}.
  \label{eq:ta_pive}
\end{equation}

Matching to \eqref{eq:pive_sol}:
\begin{itemize}
  \item $\tilde{\mathbf{Z}}' = \bone'$,
  \item $\tilde{\mathbf{D}} = \boldsymbol{\kappa}$,
  \item $\tilde{\mathbf{Z}}'\tilde{\mathbf{D}} = \bone'\boldsymbol{\kappa}$,
  \item $\mathbf{T}^{a} = \mathrm{diag}(\boldsymbol{\kappa}_1 - \boldsymbol{\kappa}_0)$.
\end{itemize}

By Proposition~\ref{prop:weights}:
\begin{equation}
  \boldsymbol{\omega}^{a\prime}
  = (\bone'\boldsymbol{\kappa})^{-1}\,\bone'\,\mathrm{diag}(\boldsymbol{\kappa}_1 - \boldsymbol{\kappa}_0),
  \quad\text{i.e.}\quad
  \omega_i^a = \frac{\kappa_{i1} - \kappa_{i0}}{\sum_j \kappa_j}.
  \label{eq:w_a}
\end{equation}

Again, Condition~\ref{cond:smoother} holds trivially since no outcome regression is involved.

\subsection{Checking normalization via Shortcut SC1}

We apply SC1 by checking $\mathbf{T}^a\bone$:
\begin{align*}
  \mathbf{T}^a\bone
  &= \mathrm{diag}(\boldsymbol{\kappa}_1 - \boldsymbol{\kappa}_0)\,\bone
   = \boldsymbol{\kappa}_1 - \boldsymbol{\kappa}_0.
\end{align*}
This is \emph{not} the zero vector in general,
since $\kappa_{i1} - \kappa_{i0} \neq 0$ whenever $Z_i \neq \phat_i$.
Therefore SC1 fails, and $\sum_i \omega_i^a \neq 0$ in general.

\begin{keyresult}[$\that{a}$ is not normalized]
$\mathbf{T}^a\bone = \boldsymbol{\kappa}_1 - \boldsymbol{\kappa}_0 \neq \bzero$
in general.
Hence $\sum_i \omega_i^a \neq 0$, so $\that{a}$ is \textbf{not normalized}.
No additional conditions can fix this --- it is a structural feature of the estimator.
\end{keyresult}

\subsection{Checking SC2: treated weights}

We apply SC2 by checking whether $\mathbf{T}^a\mathbf{D} = \tilde{\mathbf{D}} = \boldsymbol{\kappa}$:
\begin{align*}
  \mathbf{T}^a\mathbf{D}
  &= \mathrm{diag}(\boldsymbol{\kappa}_1 - \boldsymbol{\kappa}_0)\,\mathbf{D}
   = \boldsymbol{\kappa}_1\circ\mathbf{D} - \boldsymbol{\kappa}_0\circ\mathbf{D}.
\end{align*}
Since $\kappa_{i0}D_i = 0$ for all $i$:
\[
  \mathbf{T}^a\mathbf{D} = \boldsymbol{\kappa}_1 \neq \boldsymbol{\kappa} = \tilde{\mathbf{D}}^a.
\]
SC2 fails. Together with SC1 failing, we check SC3:

\begin{align*}
  \mathbf{T}^a(\bone - \mathbf{D})
  &= \mathrm{diag}(\boldsymbol{\kappa}_1 - \boldsymbol{\kappa}_0)\,(\bone - \mathbf{D})
   = (\boldsymbol{\kappa}_1 - \boldsymbol{\kappa}_0)\circ(\bone - \mathbf{D}).
\end{align*}
Since $\kappa_{i1}(1-D_i) = 0$ for all $i$:
\[
  \mathbf{T}^a(\bone - \mathbf{D}) = -\boldsymbol{\kappa}_0.
\]
Is $-\boldsymbol{\kappa}_0 = -\tilde{\mathbf{D}}^a = -\boldsymbol{\kappa}$?
No, since $\kappa_i \neq \kappa_{i0}$ in general.
SC3 also fails, confirming that $\that{a}$ is \textbf{fully unnormalized}.

\begin{keyresult}[$\that{a}$ is fully unnormalized]
All three shortcut checks fail:
\[
  \sum_i \omega_i^a \neq 0,\quad
  \sum_i \omega_i^a D_i \neq 1,\quad
  \sum_i \omega_i^a (1-D_i) \neq -1 \quad\text{in general}.
\]
This confirms the classification in Knaus (2024, Appendix A.4) and SUW (2024).
The practical implication is that $\that{a}$ is not translation invariant:
adding a constant $c$ to every $Y_i$ changes the estimate.
\end{keyresult}

% ===============================================================
\section{Derivation: \texorpdfstring{$\hat\tau_u$}{tau\_u} (Uysal Normalized)}
\label{sec:tau_u}
% ===============================================================

\subsection{Where \texorpdfstring{$\hat\tau_u$}{tau\_u} sits in Knaus (2024)}

Knaus (2024, Appendix~A.4) identifies $\that{u}$ as the Wald-IPW estimator
\emph{with} Condition~\ref{cond:ipwnorm}(C6b) applied to the instrument-group
inverse probability weights.
In Knaus's Table~5, Wald-IPW with C6b is fully normalized.
Our task here is to make this completely explicit by deriving the transformation
matrix $\mathbf{T}^u$, the PIVE ingredients, and verifying normalization
step by step using the same matrix machinery as for $\that{a,10}$.

\subsection{Step 1: Define the normalized instrument IPW weights}
\label{ssec:u_step1}

The raw instrument inverse probability weights are
\begin{equation}
  \lambda_{1,i}^{\mathrm{ipw},z} = \frac{Z_i}{\phat_i},
  \qquad
  \lambda_{0,i}^{\mathrm{ipw},z} = \frac{1-Z_i}{1-\phat_i}.
  \label{eq:raw_ipw_z}
\end{equation}
In vector notation: $\boldsymbol{\lambda}_1^z = \mathrm{diag}(Z_i/\phat_i)_{i=1}^N$
acting on $\bone$, i.e.\ $\boldsymbol{\lambda}_1^z = (\lambda_{1,1}^{\mathrm{ipw},z},\ldots,\lambda_{1,N}^{\mathrm{ipw},z})'$,
and analogously for $\boldsymbol{\lambda}_0^z$.

\medskip\noindent
\textbf{Condition 6b} (Knaus 2024, Condition~6) applied to instrument groups requires:
\begin{equation}
  \lambda_{1,i}^{\mathrm{norm},z}
    := \frac{\lambda_{1,i}^{\mathrm{ipw},z}}{\EN[\lambda_1^{\mathrm{ipw},z}]}
     = \frac{Z_i/\phat_i}{\frac{1}{N}\sum_j Z_j/\phat_j}
     = \frac{N\,Z_i/\phat_i}{s_1},
  \quad
  \lambda_{0,i}^{\mathrm{norm},z}
    = \frac{N\,(1-Z_i)/(1-\phat_i)}{s_0},
  \label{eq:norm_ipw_z}
\end{equation}
where we define
\begin{equation}
  s_1 := \bone'\boldsymbol{\lambda}_1^z = \sum_{i=1}^N \frac{Z_i}{\phat_i},
  \qquad
  s_0 := \bone'\boldsymbol{\lambda}_0^z = \sum_{i=1}^N \frac{1-Z_i}{1-\phat_i}.
  \label{eq:s1s0}
\end{equation}
By construction of \eqref{eq:norm_ipw_z}:
\begin{equation}
  \bone'\boldsymbol{\lambda}_1^{\mathrm{norm},z}
  = \sum_i \frac{N\,Z_i/\phat_i}{s_1}
  = \frac{N}{s_1}\sum_i \frac{Z_i}{\phat_i}
  = \frac{N \cdot s_1}{s_1} = N,
  \label{eq:c6b_sum}
\end{equation}
and identically $\bone'\boldsymbol{\lambda}_0^{\mathrm{norm},z} = N$.
This is exactly what Condition~\ref{cond:ipwnorm}(C6b) requires.

\subsection{Step 2: Write \texorpdfstring{$\hat\tau_u$}{tau\_u} in Wald-IPW form}
\label{ssec:u_step2}

The Uysal estimator is a Wald ratio: the (normalized IPW) reduced form
divided by the (normalized IPW) first stage.
Using the normalized weights \eqref{eq:norm_ipw_z}, define:
\begin{align}
  \text{Reduced form numerator:}&\quad
    N^{-1}\bone'\,\mathrm{diag}(\boldsymbol{\lambda}_1^{\mathrm{norm},z} - \boldsymbol{\lambda}_0^{\mathrm{norm},z})\,\mathbf{Y},
  \label{eq:rf_num}\\[4pt]
  \text{First stage denominator:}&\quad
    N^{-1}\bone'\,\mathrm{diag}(\boldsymbol{\lambda}_1^{\mathrm{norm},z} - \boldsymbol{\lambda}_0^{\mathrm{norm},z})\,\mathbf{D}.
  \label{eq:fs_denom}
\end{align}
Substituting the definition \eqref{eq:norm_ipw_z}, the reduced form numerator
\eqref{eq:rf_num} becomes:
\begin{align*}
  N^{-1}\sum_i (\lambda_{1,i}^{\mathrm{norm},z} - \lambda_{0,i}^{\mathrm{norm},z})\,Y_i
  &= N^{-1}\sum_i \left(\frac{N Z_i/\phat_i}{s_1} - \frac{N(1-Z_i)/(1-\phat_i)}{s_0}\right) Y_i\\
  &= \frac{\sum_i Z_i Y_i/\phat_i}{s_1} - \frac{\sum_i (1-Z_i)Y_i/(1-\phat_i)}{s_0}\\
  &= \mu_{Y,1}^{\mathrm{norm}} - \mu_{Y,0}^{\mathrm{norm}},
\end{align*}
and the first stage denominator \eqref{eq:fs_denom} analogously yields
$\mu_{D,1}^{\mathrm{norm}} - \mu_{D,0}^{\mathrm{norm}} =: \Delta_D$.
Hence
\begin{equation}
  \that{u}
  = \frac{\mu_{Y,1}^{\mathrm{norm}} - \mu_{Y,0}^{\mathrm{norm}}}{\mu_{D,1}^{\mathrm{norm}} - \mu_{D,0}^{\mathrm{norm}}}
  = \frac{N^{-1}\bone'\,\mathrm{diag}(\boldsymbol{\lambda}_1^{\mathrm{norm},z} - \boldsymbol{\lambda}_0^{\mathrm{norm},z})\,\mathbf{Y}}%
         {N^{-1}\bone'\,\mathrm{diag}(\boldsymbol{\lambda}_1^{\mathrm{norm},z} - \boldsymbol{\lambda}_0^{\mathrm{norm},z})\,\mathbf{D}}.
  \label{eq:tau_u_matrix}
\end{equation}

\subsection{Step 3: PIVE identification}
\label{ssec:u_step3}

Matching \eqref{eq:tau_u_matrix} to the PIVE solution
$\hat\tau = (\tilde{\mathbf{Z}}'\tilde{\mathbf{D}})^{-1}\tilde{\mathbf{Z}}'\tilde{\mathbf{Y}}$,
we read off:

\begin{center}
\renewcommand{\arraystretch}{1.5}
\begin{tabular}{@{}lll@{}}
\toprule
PIVE component & Expression & Dimension \\
\midrule
Pseudo-instrument & $\tilde{\mathbf{Z}}' = N^{-1}\bone'$ & $1\times N$ \\
Pseudo-treatment  & $\tilde{\mathbf{D}} = \mathrm{diag}(\boldsymbol{\lambda}_1^{\mathrm{norm},z} - \boldsymbol{\lambda}_0^{\mathrm{norm},z})\,\bone$ & $N\times 1$ \\
Transformation matrix & $\mathbf{T}^u = \mathrm{diag}(\boldsymbol{\lambda}_1^{\mathrm{norm},z} - \boldsymbol{\lambda}_0^{\mathrm{norm},z})$ & $N\times N$ \\
Normalizing scalar & $\tilde{\mathbf{Z}}'\tilde{\mathbf{D}} = N^{-1}\bone'\mathbf{T}^u\bone = \Delta_D$ & scalar \\
\bottomrule
\end{tabular}
\end{center}

\medskip\noindent
To see that the normalizing scalar equals $\Delta_D$, note:
\[
  N^{-1}\bone'\mathbf{T}^u\bone
  = N^{-1}\sum_i (\lambda_{1,i}^{\mathrm{norm},z} - \lambda_{0,i}^{\mathrm{norm},z})
  = N^{-1}\bone'\,\mathrm{diag}(\boldsymbol{\lambda}_1^{\mathrm{norm},z}-\boldsymbol{\lambda}_0^{\mathrm{norm},z})\,\mathbf{D}
  = \Delta_D.
\]
(The second equality holds because we define $\Delta_D$ as the first-stage,
i.e.\ the same expression evaluated with $\mathbf{D}$ instead of $\bone$. Wait ---
more precisely, $\tilde{\mathbf{Z}}'\tilde{\mathbf{D}} = N^{-1}\bone'\tilde{\mathbf{D}}$,
and since $\tilde{D}_i = (\lambda_{1,i}^{\mathrm{norm},z} - \lambda_{0,i}^{\mathrm{norm},z})D_i$
is element $i$ of $\mathbf{T}^u\mathbf{D}$, we have
$\tilde{\mathbf{Z}}'\tilde{\mathbf{D}} = N^{-1}\bone'\mathbf{T}^u\mathbf{D} = \Delta_D$.)

\medskip\noindent
Since the kappa weights involve no outcome regression, Condition~\ref{cond:smoother}
(smoother matrix) holds trivially. By Proposition~\ref{prop:weights}, the
outcome weights are:
\begin{equation}
  \boldsymbol{\omega}^{u\prime}
  = (\tilde{\mathbf{Z}}'\tilde{\mathbf{D}})^{-1}\,\tilde{\mathbf{Z}}'\,\mathbf{T}^u
  = \Delta_D^{-1}\,N^{-1}\,\bone'\,
    \mathrm{diag}(\boldsymbol{\lambda}_1^{\mathrm{norm},z} - \boldsymbol{\lambda}_0^{\mathrm{norm},z}).
  \label{eq:w_u_matrix}
\end{equation}
Per observation:
\begin{equation}
  \boxed{
    \omega_i^u
    = \frac{\lambda_{1,i}^{\mathrm{norm},z} - \lambda_{0,i}^{\mathrm{norm},z}}{N\,\Delta_D}
    = \frac{1}{\Delta_D}
      \left(\frac{Z_i}{\phat_i\,s_1} - \frac{1-Z_i}{(1-\phat_i)\,s_0}\right).
  }
  \label{eq:w_u_scalar}
\end{equation}

\subsection{Step 4: SC1 --- do the outcome weights sum to zero?}
\label{ssec:u_sc1}

The general shortcut SC1 checks $\mathbf{T}^u\bone$.
In full matrix form:
\begin{equation}
  \mathbf{T}^u\bone
  = \mathrm{diag}(\boldsymbol{\lambda}_1^{\mathrm{norm},z} - \boldsymbol{\lambda}_0^{\mathrm{norm},z})\,\bone
  = \boldsymbol{\lambda}_1^{\mathrm{norm},z} - \boldsymbol{\lambda}_0^{\mathrm{norm},z}.
  \label{eq:tu_1N}
\end{equation}
This is not $\bzero$ in general (each element is $\lambda_{1,i}^{\mathrm{norm},z} - \lambda_{0,i}^{\mathrm{norm},z} \neq 0$
unless $Z_i = \phat_i$), so SC1 via the transformation matrix alone does not deliver the result.
We must therefore work with the \emph{full weight expression}
$\boldsymbol{\omega}^{u\prime} = \Delta_D^{-1} N^{-1}\bone'\mathbf{T}^u$ directly.

\medskip\noindent
Compute $\sum_i \omega_i^u = \boldsymbol{\omega}^{u\prime}\bone$:
\begin{align}
  \boldsymbol{\omega}^{u\prime}\bone
  &= \Delta_D^{-1}\,N^{-1}\,\bone'\,\mathbf{T}^u\,\bone
  \notag\\
  &= \Delta_D^{-1}\,N^{-1}\,\bone'\,(\boldsymbol{\lambda}_1^{\mathrm{norm},z} - \boldsymbol{\lambda}_0^{\mathrm{norm},z})
  \notag\\
  &= \Delta_D^{-1}\,N^{-1}\,
     \left(\bone'\boldsymbol{\lambda}_1^{\mathrm{norm},z} - \bone'\boldsymbol{\lambda}_0^{\mathrm{norm},z}\right).
  \label{eq:u_sc1_step1}
\end{align}
Now apply Condition~\ref{cond:ipwnorm}(C6b) as established in \eqref{eq:c6b_sum}:
\[
  \bone'\boldsymbol{\lambda}_1^{\mathrm{norm},z} = N
  \quad\text{and}\quad
  \bone'\boldsymbol{\lambda}_0^{\mathrm{norm},z} = N.
\]
Substituting into \eqref{eq:u_sc1_step1}:
\begin{equation}
  \boldsymbol{\omega}^{u\prime}\bone
  = \Delta_D^{-1}\,N^{-1}\,(N - N)
  = \Delta_D^{-1}\cdot 0 = 0. \qquad\checkmark
  \label{eq:u_norm}
\end{equation}

\begin{keyresult}[$\that{u}$ is normalized --- and why C6b is the key]
$\sum_i \omega_i^u = 0$ holds \emph{if and only if}
$\bone'\boldsymbol{\lambda}_1^{\mathrm{norm},z} = \bone'\boldsymbol{\lambda}_0^{\mathrm{norm},z}$,
i.e.\ exactly Condition~\ref{cond:ipwnorm}(C6b).
Without C6b (i.e.\ using raw $\lambda_1^{\mathrm{ipw},z}$ instead of $\lambda_1^{\mathrm{norm},z}$),
the two sums $\bone'\boldsymbol{\lambda}_1^z = s_1$ and $\bone'\boldsymbol{\lambda}_0^z = s_0$
need not be equal, so $\boldsymbol{\omega}^{u\prime}\bone = \Delta_D^{-1}N^{-1}(s_1 - s_0) \neq 0$
in general.
This is precisely why $\that{a,1} = \that{t}$ (Tan/Frölich, which uses raw IPW weights)
is \emph{not} normalized, while $\that{u}$ (Uysal, which applies C6b) is.
\end{keyresult}

\subsection{Step 5: SC2 --- do the treated outcome weights sum to one?}
\label{ssec:u_sc2}

We compute $\boldsymbol{\omega}^{u\prime}\mathbf{D} = \sum_i \omega_i^u D_i$:
\begin{align}
  \boldsymbol{\omega}^{u\prime}\mathbf{D}
  &= \Delta_D^{-1}\,N^{-1}\,\bone'\,\mathbf{T}^u\,\mathbf{D}
  \notag\\
  &= \Delta_D^{-1}\,N^{-1}\,\bone'\,
     \mathrm{diag}(\boldsymbol{\lambda}_1^{\mathrm{norm},z} - \boldsymbol{\lambda}_0^{\mathrm{norm},z})\,\mathbf{D}.
  \label{eq:u_sc2_step1}
\end{align}
Element-wise, the product inside is
$(\lambda_{1,i}^{\mathrm{norm},z} - \lambda_{0,i}^{\mathrm{norm},z})\,D_i$, so:
\begin{align}
  \bone'\,\mathrm{diag}(\boldsymbol{\lambda}_1^{\mathrm{norm},z} - \boldsymbol{\lambda}_0^{\mathrm{norm},z})\,\mathbf{D}
  &= \sum_i (\lambda_{1,i}^{\mathrm{norm},z} - \lambda_{0,i}^{\mathrm{norm},z})\,D_i
  \notag\\
  &= \sum_i \left(\frac{N Z_i D_i/\phat_i}{s_1} - \frac{N(1-Z_i)D_i/(1-\phat_i)}{s_0}\right)
  \notag\\
  &= N\left(\frac{\sum_i Z_i D_i/\phat_i}{s_1} - \frac{\sum_i (1-Z_i)D_i/(1-\phat_i)}{s_0}\right)
  \notag\\
  &= N\left(\mu_{D,1}^{\mathrm{norm}} - \mu_{D,0}^{\mathrm{norm}}\right)
   = N\,\Delta_D.
  \label{eq:u_sc2_step2}
\end{align}
In the last step we use the definition $\Delta_D = \mu_{D,1}^{\mathrm{norm}} - \mu_{D,0}^{\mathrm{norm}}$.
Substituting \eqref{eq:u_sc2_step2} back into \eqref{eq:u_sc2_step1}:
\begin{equation}
  \boldsymbol{\omega}^{u\prime}\mathbf{D}
  = \Delta_D^{-1}\,N^{-1}\,N\,\Delta_D
  = 1. \qquad\checkmark
  \label{eq:u_treated_sum}
\end{equation}

\medskip\noindent
Note what happened here: the same quantity $\Delta_D$ appears in both numerator
and denominator and cancels exactly.
This is not a coincidence --- it is a consequence of defining
$\tilde{\mathbf{Z}}'\tilde{\mathbf{D}} = N^{-1}\bone'\mathbf{T}^u\mathbf{D}$
as the PIVE normalizing scalar.
In other words, the denominator of $\that{u}$ is precisely built to make the
treated weights sum to one.

\subsection{Step 6: Untreated weights sum to minus one (by subtraction)}
\label{ssec:u_untreated}

Since $\sum_i \omega_i^u = 0$ (Step~4) and $\sum_i \omega_i^u D_i = 1$ (Step~5):
\begin{equation}
  \boldsymbol{\omega}^{u\prime}(\bone - \mathbf{D})
  = \boldsymbol{\omega}^{u\prime}\bone - \boldsymbol{\omega}^{u\prime}\mathbf{D}
  = 0 - 1 = -1. \qquad\checkmark
  \label{eq:u_untreated_sum}
\end{equation}

\begin{keyresult}[$\that{u}$ is fully normalized]
The outcome weights $\omega_i^u$ in \eqref{eq:w_u_scalar} satisfy:
\[
  \sum_{i=1}^N \omega_i^u = 0,\qquad
  \sum_{i=1}^N \omega_i^u D_i = 1,\qquad
  \sum_{i=1}^N \omega_i^u (1-D_i) = -1.
\]
\textbf{Which condition is responsible?}
Condition~\ref{cond:ipwnorm}(C6b) from Knaus (2024): it forces
$\bone'\boldsymbol{\lambda}_1^{\mathrm{norm},z} = \bone'\boldsymbol{\lambda}_0^{\mathrm{norm},z} = N$,
which is the single fact that drives the cancellation in Step~4.
Step~5 requires no further condition --- it follows automatically because
$\Delta_D$ is the first-stage denominator of $\that{u}$ and therefore
appears in both numerator and denominator of $\boldsymbol{\omega}^{u\prime}\mathbf{D}$.
\end{keyresult}

\subsection{What goes wrong without C6b: contrast with \texorpdfstring{$\hat\tau_{a,1}$}{tau\_a,1}}
\label{ssec:u_vs_a1}

The estimator $\that{a,1} = \that{t}$ (Tan/Frölich) uses raw IPW weights
$\lambda_{1,i}^{\mathrm{ipw},z}$ and $\lambda_{0,i}^{\mathrm{ipw},z}$ in place of
the normalized $\lambda_{1,i}^{\mathrm{norm},z}$ and $\lambda_{0,i}^{\mathrm{norm},z}$.
Its transformation matrix is therefore
\[
  \mathbf{T}^{a,1} = \mathrm{diag}(\boldsymbol{\lambda}_1^z - \boldsymbol{\lambda}_0^z),
  \quad\text{and}\quad
  \boldsymbol{\omega}^{a,1\prime}
  = \Delta_D^{a,1\,-1}\,N^{-1}\,\bone'\,\mathrm{diag}(\boldsymbol{\lambda}_1^z - \boldsymbol{\lambda}_0^z).
\]
Repeating Step~4 for $\that{a,1}$:
\begin{align*}
  \boldsymbol{\omega}^{a,1\prime}\bone
  &= \Delta_D^{a,1\,-1}\,N^{-1}\,(\bone'\boldsymbol{\lambda}_1^z - \bone'\boldsymbol{\lambda}_0^z)\\
  &= \Delta_D^{a,1\,-1}\,N^{-1}\,(s_1 - s_0).
\end{align*}
Since $s_1 = \sum_i Z_i/\phat_i$ and $s_0 = \sum_i (1-Z_i)/(1-\phat_i)$
are in general \emph{not equal}, the sum of weights is
\[
  \boldsymbol{\omega}^{a,1\prime}\bone \neq 0 \quad\text{in general},
\]
so $\that{a,1}$ is not normalized.
The failure is exactly that C6b is not applied: without rescaling by $s_1$ and $s_0$,
the two groups contribute different total weight to the estimator.
Adding a constant $c$ to every $Y_i$ shifts the estimate by $c\cdot(s_1-s_0)/(\Delta_D^{a,1}\cdot N)$,
making $\that{a,1}$ not translation invariant.

% ===============================================================
\section{Summary and Comparison}
\label{sec:summary}
% ===============================================================

Table~\ref{tab:summary} collects the results for all three estimators.

\begin{table}[h!]
\centering
\caption{Outcome weights: existence, form, and normalization}
\label{tab:summary}
\renewcommand{\arraystretch}{1.4}
\begin{tabular}{@{}lllll@{}}
\toprule
Estimator & Closed-form & $\sum_i\omega_i$ & $\sum_i\omega_iD_i$ & $\sum_i\omega_i(1-D_i)$ \\
\midrule
$\that{a}$    & $\checkmark$ (unconditional) & $\neq 0$ & $\neq 1$ & $\neq -1$ \\
$\that{a,0}$  & $\checkmark$ (unconditional) & $\neq 0$ & $\neq 1$ & $= -1$    \\
$\that{a,10}$ & $\checkmark$ (unconditional) & $= 0$    & $= 1$    & $= -1$    \\
$\that{u}$    & $\checkmark$ (unconditional) & $= 0$    & $= 1$    & $= -1$    \\
\bottomrule
\end{tabular}\\[4pt]
\small\textit{Notes:} $\that{a,1}$ (Tan/Frölich estimator) equals $\that{a}$ rescaled and is also fully unnormalized.
Normalization classification follows Knaus (2024), Table~4.
\end{table}

\medskip
\noindent\textbf{Why $\that{a,10}$ and $\that{u}$ achieve full normalization by different routes:}
\begin{itemize}
  \item $\that{a,10}$ achieves full normalization by \emph{separately} normalizing the $\kappa_1$
        and $\kappa_0$ weighted sums.
        Each half sums to 1 and -1 respectively by construction, and the cancellation is exact.
        No condition from Knaus (2024) needs to be invoked --- normalization is built in algebraically.
  \item $\that{u}$ achieves full normalization via Hájek normalization of the IPW weights,
        corresponding to Condition~\ref{cond:ipwnorm}(C6b).
        The analogy in Knaus (2024) is Wald-IPW with C6b applied (Table~5, Knaus 2024).
  \item $\that{a}$ and $\that{a,1}$ use a single shared denominator that does not separately
        normalize the treated and untreated contributions, so neither the sum-to-zero
        nor the group-level conditions are met.
\end{itemize}

% ---------------------------------------------------------------
\section*{References}
% ---------------------------------------------------------------

\begin{description}
  \item[Abadie (2003)] Semiparametric instrumental variable estimation of treatment response models.
        \textit{Journal of Econometrics}, 113(2), 231--263.

  \item[Knaus (2024)] Treatment effect estimators as weighted outcomes.
        arXiv:2411.11559v2. University of Tübingen.

  \item[Słoczy\'nski, Uysal \& Wooldridge (2024)] Abadie's kappa and weighting estimators
        of the local average treatment effect.
        \textit{Journal of Business \& Economic Statistics}, 43(1), 164--177.

  \item[Uysal (2011)] Three essays on doubly robust estimation methods.
        PhD Dissertation, University of Konstanz.

  \item[Hájek (1971)] Comment on ``An essay on the logical foundations of survey sampling, part one.''
        In \textit{Foundations of Statistical Inference}, p.~236. Toronto: Holt, Rinehart and Winston.
\end{description}

\end{document}
